\section{本章小结}
\label{sec:ospec-summary}

矩阵的特征值在算子上的对应物是谱，由$T-\lambda I$能否在整个空间上稳定求逆来定义。谱总是有界闭集，含于半径为$\lVert T\rVert$的圆盘，参数越靠近谱，反解的误差放大越严重。有限维中谱只有特征值，无限维多出两类点：连续谱上$T-\lambda I$是单射、像稠密而不满，方程到处可以近似求解，解却不稳定；剩余谱上像不稠密，有一块右端到达不了。

乘法算子把“连续的对角矩阵”这个类比说得最清楚：对角元铺满$[0,1]$，谱是整段区间，却没有特征函数，单位高台只是近似特征向量，极限不在$L^2$中。右移算子的谱是整个闭单位圆盘，内部是剩余谱，边界是连续谱，而它的有限截断只有特征值$0$，说明矩阵截断不一定逼近谱。

紧性使非零谱回到矩阵的情形：非零谱点都是特征值，特征空间有限维，特征值只能向零聚集，单射已经推出可逆。零在无限维中总是谱点，却未必是特征值。累积积分是紧算子，谱只有$\{0\}$，没有一个特征向量，所以紧性本身不保证特征展开。

在Hilbert空间中加入自伴条件，特征值是实数，特征方向两两正交，并与零空间一起张满整个空间，算子作用写成按Hilbert范数收敛的级数，截去前$m$项的算子范数误差是第一个被舍弃的特征值的绝对值。一般紧算子借助$T^*T$得到奇异分解，不要求自伴，也不要求输入输出空间相同，截断误差$\sigma_{m+1}$是秩$m$算子所能达到的最小值。

平方可积核把这些展开变成核函数的展开，核的$L^2$范数对应Frobenius范数，其平方等于特征值或奇异值的平方和；自伴核与正规核有正交特征函数展开，一般核用左右奇异函数展开。同一截断有两种误差：算子误差只看第一个舍弃项，核误差累积全部舍弃项。展开在乘积$L^2$意义下收敛，逐点或一致收敛需要另外的条件。

小奇异值使正向逼近容易、反向恢复困难，这些展开并不自行保证由输出稳定地恢复输入。算子方程的可解条件与稳定性留到第~\ref{chap:linear-operator-equations}章。
